\section{Qualche progetto}

% \experienceTitle{Adattamento online di robot}{2022}
% \experienceDetails{Adattamento continuo di un sistema robotico ad alcuni compiti e ambienti:\\\hspace*{9pt}(1) evitare le collisioni, (2) evitare le aree, (3) navigare un labirinto.}{}

\experienceTitle{Strumento per la generazione automatica di software di controllo per robot}{2025}
\experienceDetails{Creazione di un tool per la generazione automatica di software di controllo\\\textnormal{Sviluppato a IRIDIA, Bruxelles}}{ \href{https://github.com/p-baldini/AutoMoDe/tree/lahmacun}{link al repository github}}

\experienceTitle{Simulatore di dispositivo memristivo}{2021}
\experienceDetails{Creazione del simulatore di un dispositivo memristivo nanometrico\\\textnormal{per un gruppo di ricerca in fisica di Torino}}{ \href{https://p-baldini.github.io/nns.c/}{link al sito web}}

\pagebreak

\experienceTitle{Servizi di magazzino distribuito}{2021}
\experienceDetails{Sistemi eterogenei distribuiti che utilizzano il framework Jade e il linguaggio Jason\\\textnormal{Progetto di corso}}{ \href{https://github.com/p-baldini/AMW}{link al repository github}}

\experienceTitle{Strumento di induzione di regole}{2021}
\experienceDetails{Implementazione dell'algoritmo di apprendimento CN2 per l'induzione di regole\\\textnormal{Progetto di Corso}}{ \href{https://github.com/p-baldini/CN2-Algorithm}{link al repository github}}

\experienceTitle{Sistema di stima delle traiettorie}{2020}
\experienceDetails{Script Matlab per la stima della traiettoria di una palla da immagini acquisite\\\textnormal{Progetto di corso}}{ \href{https://github.com/p-baldini/TrajectoryEstimator}{link al repository github}}

\experienceTitle{Sistema di ricostruzione di traiettorie}{2020}
\experienceDetails{Script Matlab per la ricostruzione della traiettoria di una telecamera a partire da immagini acquisite\\\textnormal{Progetto di corso}}{ \href{https://github.com/p-baldini/CameraTrajectory}{link al repository github}}

\experienceTitle{Servizio/Gioco puzzle remoto e distribuito}{2020}
\experienceDetails{Servizio \& applicazione per il completamento di puzzles in maniera distribuita e cooperativa\\\textnormal{Progetto di corso}}{ \href{https://github.com/p-baldini/CooperativePuzzleService}{link al repository github}}

\experienceTitle{Web Kanban}{2020}
\experienceDetails{Sito web Kanban per la gestione collaborativa e agile dei progetti\\\textnormal{Progetto di corso}}{ \href{https://github.com/p-baldini/RutWorking}{link al repository github}}

\experienceTitle{Simulatore di particelle}{2020}
\experienceDetails{Simulatore del moto fisico delle particelle.}{}

\experienceTitle{Compilatore per un linguaggio di programmazione personalizzato}{2020}
\experienceDetails{Progettazione di un compilatore per un linguaggio di programmazione funzionale e orientato agli oggetti, per e con la collaborazione di un professore universitario.}{}

\experienceTitle{Sistema IoT: creazione hardware e software}{2019}
\experienceDetails{Collaborazione alla creazione dell'hardware di dispositivi IoT, come la connessione di sensori e la calibrazione della potenza, con la collaborazione di un professore universitario. Creazione del software di controllo per la loro interazione.}{}

\experienceTitle{CC-Tan game}{2018}
\experienceDetails{Versione Java del videogioco Android CC-Tan.}{}
